\section{模型评价与结论}
本文围绕三个递进目标，分别把节点顺序、物理空间和流水等待显式化。生命周期收缩在固定操作顺序下不增峰值；连续分配通过完整驻留状态与有代价的恢复机制获得可行方案；固定驻留图重排保持地址先后和搬运事件不变，从而将时间优化转化为受资源约束的拓扑排序。

六图中，问题一只有卷积小图相对收缩后的编号基准降低峰值，降幅为 7.62\%。问题二的搬运量依次为 44766、75546、3620、32852、34944 和 460800。问题三在这些搬运量均不增加的条件下，将总时间降低 1.30\% 至 33.38\%；其中注意力大图从 225902 降至 150499 周期。阶段比较表明，四图的改善来自固定驻留重排，两个矩阵乘图还需要先改变地址分配候选。

这些结果支持“缓存方案确定后仍存在可利用的流水优化空间”，但没有证明当前搬运量或峰值最优。卷积小图接近固定驻留下界，进一步工作应考虑改变驻留关系；卷积大图与该下界仍相差 24.95\%，值得扩展单元顺序搜索。矩阵乘大图则受到输入搬运工作量的明显限制，单纯在固定事件集合内重排的空间较小。后续改进方向应由这些不同证据决定，而不是对所有图统一增加启发式复杂度。
